% F -- Three days on one clock: two with no movement at all, one with three.
\begin{figure}[!htbp]
\centering
\fitwidth{%
\begin{tikzpicture}[x=0.575cm,y=1cm,font=\small,
  rowlab/.style={anchor=south west,font=\footnotesize\bfseries,text=clmodel},
  cnt/.style={anchor=west,font=\scriptsize,text=clrule},
  nm/.style={font=\tiny\ttfamily,text=clmove,anchor=north},
]
% ---------- row 1: 2G, never moves --------------------------------------
\node[rowlab] at (0,2.88) {Stays home all day, phone on 2G};
\draw[clrule!70] (0,2.65) -- (24,2.65);
\foreach \h in {1,5,9,13,17,21}{ \fill[clbase] (\h,2.65) circle (2.2pt); }
\node[cnt] at (24.4,2.65) {6 records, all \cid{BKVCHE1}};

% ---------- row 2: 3G, same non-day -------------------------------------
\node[rowlab] at (0,1.83) {The very same day, phone on 3G};
\draw[clrule!70] (0,1.60) -- (24,1.60);
\foreach \h in {0.5,2.5,4.5,6.5,8.5,10.5,12.5,14.5,16.5,18.5,20.5,22.5}{
  \fill[clbase] (\h,1.60) circle (2.2pt); }
\node[cnt] at (24.4,1.60) {12 records, all \cid{UKVCHE102}};

% ---------- row 3: someone who commutes ---------------------------------
\node[rowlab] at (0,0.78) {Commutes to work, and comes home};
\draw[clrule!70] (0,0.55) -- (24,0.55);
\foreach \h in {1,5,11,14,17,23}{ \fill[clbase] (\h,0.55) circle (2.2pt); }
\foreach \h in {8,18.5,21}{ \fill[clmove] (\h,0.55) circle (3.0pt); }
\node[nm] at (8,0.30)    {BKVPAN1};
\node[nm] at (18.2,0.30) {BSOCHO1};
\node[nm] at (21.6,0.30) {BKVCHE1};
\node[cnt] at (24.4,0.55) {9 records, 3 of them a move};

% ---------- the clock ----------------------------------------------------
\draw[clrule] (0,0.05) -- (24,0.05);
\foreach \h in {0,4,8,12,16,20,24}{
  \draw[clrule] (\h,0.05) -- (\h,-0.09);
  \node[anchor=north,font=\scriptsize,text=clrule] at (\h,-0.13) {\h h};
}

% ---------- the reading --------------------------------------------------
\node[draw=clrule,fill=clsoft,rounded corners=3pt,align=left,inner sep=7pt,
      text width=15.0cm,font=\footnotesize,anchor=north west] at (0,-0.72)
  {\textbf{Neither of the first two rows contains a single movement.} A mast
   asks every phone attached to it to announce itself on a fixed cycle, so an
   untouched phone still writes records all day. How many somebody has is
   therefore set partly by which radio generation their phone attached to,
   2G every four hours and 3G every two, and not by what they did. Only the
   orange dots are a change of antenna, and only there can a model beat
   \emph{``the same antenna as a moment ago''}.};
\end{tikzpicture}}
\caption{A record is not a movement}
\label{fig:record-not-movement}
\figtext{%
\figpara{Network ping frequency versus spatial movement in Figure~\ref{fig:record-not-movement}.}{Each horizontal line in Figure~\ref{fig:record-not-movement} is one person's day from midnight to midnight, and one dot is one record
written by the network. Blue dots carry the same antenna name as the dot before
them, so nothing happened; orange dots carry a different one. The first two rows of Figure~\ref{fig:record-not-movement}
are the same behaviour, somebody at home all day, seen through two generations of
radio, which is why the record count of a person says as much about their phone as
about them. And a long run of one antenna name is what makes ``repeat the last
name you saw'' such a hard competitor.}}

\end{figure}
